\chapter{The Bühlmann--Straub special case}\label{ch:bs}
% Thesis §2.6, eqs. (28)-(30). Lean: Credibility/BuhlmannStraub.lean

\begin{definition}[Bühlmann--Straub assumptions]\label{def:bsModel}
  \lean{Credibility.bsModel}
  \leanok
  \uses{def:Model}
  $\sigma_{jt}^2 = \sigma^2/w_{jt}$ with volumes $w_{jt} > 0$, and
  $\rho_{ij} = \delta_{ij}$ \eqth{28}. Write $w_{j*} = \sum_t w_{jt}$.
\end{definition}

\begin{lemma}[Volume weights]\label{lem:bs_a}
  \lean{Credibility.bs_a}
  \leanok
  \uses{def:bsModel, def:sigma2}
  $a_{jt} = w_{jt}/w_{j*}$ \eqth{29}.
\end{lemma}
\begin{proof}
  \leanok
  $\sigma_{jt}^{-2} = w_{jt}/\sigma^2$; the factor $\sigma^{-2}$ cancels.
\end{proof}

\begin{lemma}\label{lem:bs_σ_eq}
  \lean{Credibility.bs_σ_eq}
  \leanok
  \uses{def:bsModel}
  $\sigma_{jt} = \sigma/\sqrt{w_{jt}}$.
\end{lemma}
\begin{proof}
  \leanok
  $\sqrt{\sigma^2/w_{jt}} = \sqrt{\sigma^2}/\sqrt{w_{jt}}$.
\end{proof}

\begin{lemma}\label{lem:bs_σ2}
  \lean{Credibility.bs_σ2}
  \leanok
  \uses{def:bsModel, def:sigma2}
  If $T \ge 1$, $\sigma_{jj}^2 = \sigma^2/w_{j*}$ \eqth{30}.
\end{lemma}
\begin{proof}
  \leanok
  \uses{thm:U_balance, lem:bs_σ_eq, lem:U_full}
  The case $T_j = T$ of Theorem~\ref{thm:U_balance}. Directly:
  $\sigma_{jj}^2 = \sum_t a_{jt}^2\sigma_{jt}^2$ with inverse-variance weights,
  so Lemma~\ref{lem:invVarWeight_value} gives
  $\bigl(\sum_t \sigma_{jt}^{-2}\bigr)^{-1} = \bigl(\sum_t w_{jt}/\sigma^2\bigr)^{-1}
  = \sigma^2/w_{j*}$. (Equivalently, with Lemma~\ref{lem:bs_a}:
  $\sum_t (w_{jt}/w_{j*})^2\,\sigma^2/w_{jt} = \sigma^2 w_{j*}/w_{j*}^2$.)
\end{proof}

\begin{lemma}[$K$ is diagonal]\label{lem:bs_K}
  \lean{Credibility.bs_K}
  \leanok
  \uses{lem:bs_σ2, lem:K_apply}
  If $T \ge 1$, $K = \operatorname{diag}(\tau^2 + \sigma^2/w_{j*})$.
\end{lemma}
\begin{proof}
  \leanok
  \uses{lem:bs_σ2, thm:U_diag, lem:U_full}
  Theorem~\ref{thm:U_diag} with $T_j = T$ and $\rho_{ij} = \delta_{ij}$, and
  Lemma~\ref{lem:bs_σ2}.
\end{proof}

\begin{definition}[Classical credibility factor]\label{def:bsZ}
  \lean{Credibility.bsZ}
  \leanok
  $z_j^B = w_{j*}/(w_{j*} + \sigma^2/\tau^2)$, so that
  $K_{jj} = \tau^2/z_j^B$.
\end{definition}

\begin{lemma}\label{lem:bs_Kinv_one}
  \lean{Credibility.bs_Kinv_one}
  \leanok
  \uses{lem:bs_K, def:bsZ}
  If $T \ge 1$, $K^{-1}u = (z_j^B/\tau^2)_j$.
\end{lemma}
\begin{proof}
  \leanok
  \uses{thm:Core_diag}
  $K$ is diagonal (Lemma~\ref{lem:bs_K}; in Lean this is the hypothesis),
  so Theorem~\ref{thm:Core_diag}:
  $K^{-1} = \operatorname{diag}\bigl(1/(\tau^2 + \sigma^2/w_{j*})\bigr)
  = \operatorname{diag}(z_j^B/\tau^2)$.
\end{proof}

\begin{lemma}\label{lem:bs_b}
  \lean{Credibility.bs_b}
  \leanok
  \uses{lem:bs_Kinv_one, def:Model.b}
  If $T \ge 1$, $b_j = z_j^B/z_*^B$ where $z_*^B = \sum_i z_i^B$.
\end{lemma}
\begin{proof}
  \leanok
  \uses{thm:U_balance, lem:bs_σ_eq, lem:U_full}
  The case $T_j = T$ of Theorem~\ref{thm:U_balance}. Directly: by
  Lemma~\ref{lem:bs_Kinv_one}, $K^{-1}_{j*} = z_j^B/\tau^2$
  and $K^{-1}_{**} = z_*^B/\tau^2$.
\end{proof}

\begin{theorem}[Reduction to Bühlmann--Straub]\label{thm:bs_z}
  \lean{Credibility.bs_z}
  \leanok
  \uses{def:bsModel, def:bsZ, def:moments}
  If $T \ge 1$ and $J \ge 2$, the proposed factor is $z_j = z_j^B$.
\end{theorem}
\begin{proof}
  \leanok
  \uses{thm:U_balance, lem:bs_σ_eq, lem:U_full}
  The case $T_j = T$ of Theorem~\ref{thm:U_balance}. Directly:
  $K = \operatorname{diag}(\tau^2 + \sigma^2_{jj})$ (Lemmas~\ref{lem:bs_K}
  and~\ref{lem:bs_σ2}), so Theorem~\ref{thm:Core_diag} gives
  $z_j = \tau^2/(\tau^2 + \sigma^2/w_{j*}) = z_j^B$ (Lemma~\ref{lem:bs_σ2}).
  Directly: $K$ is diagonal with positive entries, hence positive definite. Then
  $s^2 = \tau^2/z_*^B = (\tau^2 + \sigma_{jj}^2)\,b_j$, so
  \[
    z_j = \frac{(1-b_j)\tau^2}{(\tau^2+\sigma_{jj}^2)(1-b_j)}
        = \frac{\tau^2}{\tau^2+\sigma_{jj}^2} = z_j^B .
  \]
  Cancelling $1 - b_j$ needs $b_j < 1$, which holds when $J \ge 2$
  because every $z_i^B > 0$.
\end{proof}

\begin{remark}[Finding 2: at least two cohorts]\label{rem:two_cohorts}
  \uses{thm:bs_z}
  With $J = 1$ we get $b = 1$ and $s^2 = \tau^2 + \sigma^2$, so
  \eqth{26} is $0/0$. The thesis never assumes $J \ge 2$, but the
  cancellation in the last step of §2.6 needs it.
\end{remark}

% ---------------------------------------------------------------------------
